\chapter{Work, Discipline, and the Productive Self}
\label{ch:work-discipline}

\epigraph{Choose a job you love, and you will never have to work a day in your life.}{Attributed to Confucius (apocryphal; popularized by motivational industry)}

\noindent
The moral prestige of overwork, the normalization of precarity, and the ideology of ``hustle culture'' are not spontaneous developments.
They are doxonomic products that serve specific institutional interests.

\section{The Moral Prestige of Overwork}
\label{sec:ch26-overwork}

\begin{definition}[Work Ethic Regime]
\label{def:work-ethic}
\index{work ethic}
The \emph{work ethic regime} is an anthropological regime $\alpha_W$ holding that:
\begin{enumerate}[nosep]
  \item work is the primary source of moral worth,
  \item leisure is suspect (``idle hands''),
  \item suffering through labor is ennobling,
  \item refusal to work is moral failure.
\end{enumerate}
\end{definition}

This regime predates capitalism (Weber's Protestant ethic) but was systematically amplified by industrial and post-industrial capital, which benefits from a labor force that accepts overwork as virtue rather than extraction \autocite{sennett1998}.

\section{Labor Discipline as Internalized Belief}
\label{sec:ch26-internalized}

\begin{model}[Internalized Discipline]
\label{mod:discipline}
\index{labor discipline}
Let $w_j$ be the wage and $h_j$ the hours worked by agent $j$.
Under external discipline, $h_j = h^*$ is enforced by monitoring.
Under doxonomic discipline, $j$ chooses $h_j > h^*$ voluntarily because the work ethic regime increases $j$'s subjective utility from labor:
\[
  u_j(h) = w_j h - c(h) + \belief{W} \cdot v(h)
\]
where $v(h)$ is the moral satisfaction from work and $\belief{W}$ is the strength of the work ethic belief.
The employer saves monitoring costs and extracts additional labor without additional compensation.
\end{model}

\section{Notes and References}
\label{sec:ch26-notes}

On the moral history of work, see \textcite{sennett1998}.
Precarity and the new labor regime are analyzed in \textcite{standing2011}.
On self-exploitation, see \textcite{brown2015}.
For the doxonomic function of management ideology, see \textcite{bourdieu1991}.

\section{Exercises}

\begin{exercise}
Estimate the monetary value to employers of the work ethic regime by computing the difference in labor supply between $\belief{W} = 0.8$ and $\belief{W} = 0.2$ using Model~\ref{mod:discipline}.
\end{exercise}

\begin{exercise}
Analyze ``hustle culture'' as a doxonomic product.
Identify the funders, institutions, and channels.
\end{exercise}

\begin{exercise}
Design an experiment to test whether exposure to hustle-culture narratives increases willingness to accept unpaid overtime.
\end{exercise}
